\begin{table}[!htbp]
\centering
\caption{Distributional statistics}
\begin{tabular}{lccccccc}
\toprule
Variable & $cv$ & Gini & $q_1$ & $q_2$ & $q_3$ & $q_4$ & $q_5$ \\
\midrule
\multicolumn{8}{l}{Hours} \\
Model $E_0$ & 0.04 & NaN & 0.40 & 0.41 & 0.42 & 0.42 & 0.44 \\
Data (CPS) & 0.22 & 0.11 & 0.24 & 0.31 & 0.33 & 0.35 & 0.42 \\
\addlinespace
\multicolumn{8}{l}{Earnings} \\
Model $E_0$ & 0.62 & 0.30 & 8.8\% & 12.6\% & 16.3\% & 23.7\% & 38.6\% \\
Data (CPS) & 0.56 & 0.29 & 7.9\% & 13.7\% & 18.0\% & 23.3\% & 37.1\% \\
Data (SCF) & 2.65 & 0.61 & -0.2\% & 4.0\% & 13.0\% & 22.9\% & 60.2\% \\
\addlinespace
\multicolumn{8}{l}{Wealth} \\
Model $E_0$ & 0.00 & -2.16 & 0.0\% & 81.0\% & 6.3\% & 6.3\% & 6.3\% \\
Data (SCF) & 6.53 & 0.80 & -0.3\% & 1.3\% & 5.0\% & 12.2\% & 81.7\% \\
\bottomrule
\end{tabular}
\\[3ex]
\raggedright\footnotesize{\textit{Notes.} $cv$ refers to coefficient of variation. $q_1, \dots, q_5$ refer, for earnings and wealth, to the share held by all people in the corresponding quintile with respect to the total. However, for hours it is the average number of hours worked by people in the corresponding quintile. Statistics from SCF correspond to the 1998 wave and are quoted from Budria et al. (2002). Statistics from CPS correspond to the 2002 wave.}\\
\normalsize
\end{table}
